\begin{comment}
The Goldreich--Goldwasser--Micali (\GGM) construction~\cite{GGM86} was originally
proposed as a way to build pseudorandom functions from any pseudorandom generator
(\PRG).  In the context of post-quantum digital signatures based on zero-knowledge
proofs, the \GGM\ tree serves a completely different but equally central purpose:
it is a \emph{seed-compression mechanism} that allows a signer to commit to $t$
independent ephemeral seeds using only a single master seed, while later revealing
a subset of those seeds to a verifier without exposing the hidden ones.

This role makes the \GGM\ tree a \emph{single point of failure} with respect to
physical attacks.  Any fault that causes the tree to misclassify a hidden seed as
publishable immediately breaks the Zero-Knowledge (ZK) property of the underlying
protocol, leaking secret-equivalent information.  The key observation of this work
is that this vulnerability is \emph{not specific to any one scheme}---it is a
structural consequence of how the \GGM\ tree is used in all ZK-based post-quantum
signature schemes, including LESS-v2, MEDS, CROSS, FAEST, Mirath, SDitH, and others.

\paragraph{Contributions.}
This paper provides the following:
\begin{enumerate}
  \item A description of the original \GGM\ construction and the punctured-PRF
        property that underlies all uses in ZK signatures (Section~\ref{sec:background}).
  \item A scheme-by-scheme analysis of how LESS-v2, MEDS, CROSS, and the
        MPCitH/VOLEitH family each instantiate the \GGM\ tree and the associated
        flag/reference-tree mechanism.
  \item A unified abstract model, the \emph{Generic ZK Seed Tree} (\GZKST),
        that captures all these instantiations
        (Section~\ref{sec:abstract}).
  \item A taxonomy of fault models applicable to the \GZKST\ abstraction and
        a Generic Leakage Theorem showing that any effective fault yields a
        dual-revelation and enables recovery of a secret-equivalent component
        (Section~\ref{sec:theoretical}).
\end{enumerate}
\end{comment}

\subsection{Related Work}
\label{sec:relwork}

In the context of fault attacks against post-quantum signature schemes, several 
works have been proposed in recent years, particularly targeting schemes that rely on 
\GGM{}-tree-like constructions or related seed-tree compression techniques, such as 
LESS~\cite{LESS-spec}, CROSS~\cite{CROSSspec}, and MPC-in-the-Head-based 
schemes~\cite{ryde2025,FAESTv22025}.

In~\cite{ZKFault}, the authors present a fault attack that exploits the tree-based 
compression mechanism used in zero-knowledge-based code-based signature schemes. 
Their primary target is the first version of LESS~\cite{LESS-spec}, for which they 
provide a detailed analysis showing how fault injections in the seed-tree construction 
can lead to the recovery of secret information and ultimately the signing key. The 
work further extends the attack analysis to CROSS and briefly 
discusses how similar ideas could be adapted to MEDS~\cite{ChouNPRRST23}, since 
these schemes rely on related seed-tree compression techniques.

More recently,~\cite{FaultToForge2025} revisits these attacks in the context of the second 
version of LESS (LESS-v2). Since LESS-v2 replaces the complete \GGM{}-tree structure of 
LESS-v1 with an incomplete tree construction, the original attack from~\cite{ZKFault} no 
longer applies directly. The authors therefore introduce new fault-assisted forgery 
attacks targeting the incomplete binary decision-tree generation and seed-publication 
process. In particular, they identify several fault-injection surfaces within the 
incomplete tree-generation logic and propose practical instruction-skip-based fault 
models capable of disclosing hidden ephemeral seeds. A key contribution of the work 
is the observation that recovering the full signing key is unnecessary for universal 
forgery; instead, recovering secret-equivalent information, namely the permutations 
associated with the secret monomial matrices, is sufficient to forge valid signatures.

Still attacking LESS-v2,~\cite{HuangHHYSTY26} presents an alternative attack strategy 
against LESS-v2 that completely avoids targeting the complex tree structure itself. 
Instead, the attack focuses on the simple \texttt{for-loop} responsible for 
transforming the digest vector $b$ into the binary vector $f$ before the incomplete 
\GGM{}-tree is generated. Using \emph{iteration-skip} and \emph{loop-abort} fault models, 
the authors exploit the linear and deterministic nature of this preprocessing stage to recover 
secret information with high probability. Unlike~\cite{FaultToForge2025}, whose main goal is 
universal forgery through recovery of secret-equivalent permutations, the attack in~\cite{HuangHHYSTY26} aims 
at full key recovery while also eliminating the need to heuristically target fault locations 
within the incomplete tree-generation procedure.

Our work unifies these scattered analyses under a single Generic ZK Seed Tree (GZKST) mode, 
formally proving the correctness and ZK-hiding invariants that all prior attacks implicitly rely on, 
and extending practical fault attacks to MEDS -- a scheme not deep previously studied in this context.

\begin{comment}
Table~\ref{tab:relwork} summarises the prior art on fault attacks against
\GGM-tree-based post-quantum signature schemes.

\begin{table}[tb]
\centering
\caption{Comparison with prior fault-attack analyses on ZK-based PQ signatures.}
\label{tab:relwork}
\setlength{\tabcolsep}{4pt}
\resizebox{\textwidth}{!}{
\begin{tabular}{@{}llllc@{}}
\toprule
Work & Target & Tree type & Surfaces & Generic? \\
\midrule
ZKFault~\cite{ZKFault}
  & LESS-v1, CROSS & complete & 3 & partial \\
Fault to Forge~\cite{FaultToForge2025}
  & LESS-v2        & incomplete & 10 & no \\
Huang et al.~\cite{HuangHHYSTY26}
  & LESS-v2        & incomplete & 1  & no \\
Jendral \& Dubrova~\cite{JendralD25}
  & FAEST (VOLEitH)& complete   & 4  & partial \\
Banda et al.~\cite{BandaBK26}
  & Mirath, RYDE+4 & forest     & key-rec.\ + forgery & MPCitH only \\
\midrule
\textbf{This work}
  & LESS-v2, MEDS, CROSS, MPCitH
  & complete + incomplete & all phases & \textbf{yes (\GZKST)} \\
\bottomrule
\end{tabular}
}
\end{table}
\end{comment}